\documentclass{article}
\usepackage[margin=1in]{geometry}
\usepackage{amsmath}
\usepackage{amssymb}
\usepackage{enumitem}
\usepackage{hyperref}

\title{Prediction And Falsification Protocols From Finite-Capacity Causal Networks}
\author{Abraham Greenman}
\date{2026}

\begin{document}
\maketitle

\begin{abstract}
This paper closes the Paper 15 prediction-and-falsification protocol theorem
as a conditional, finite, network-native interface theorem downstream of the
closed Paper 14 discriminating benchmarks certificate. The result constructs
finite protocol records, target and regime descriptors, threshold and
rejection-condition descriptors, Paper 14 compatibility rows, stability and
reproducibility descriptors, and a fail-closed no-hidden-promotion audit. The
final certificate closes only this finite protocol-structure theorem. It does
not assert protocol recovery, benchmark success, prediction success,
falsification success, physical validation, empirical adequacy, physical
promotion, physical nature realization, or a unified field theory.
\end{abstract}

\section{Scope}

Paper 15 starts from the frozen Paper 14 conditional discriminating benchmarks
certificate:
\[
\texttt{paper14\_dbm008\_final\_conditional\_certificate\_closes\_discriminating\_benchmarks\_theorem}.
\]
The Paper 14 frozen commit is:
\[
\texttt{ad4f876a1699874cd6efd7fe73d318e64f5bbe19}.
\]

The theorem closed here is intentionally local. It says that prediction and
falsification protocols can be represented as finite, auditable records in the
finite-capacity causal-network theorem ladder. It does not say that any
prediction has succeeded, any alternative has been falsified, any benchmark
has succeeded, or any physical theory has been validated.

\section{Claim Boundary}

The following claims remain false in this paper:

\begin{itemize}[leftmargin=*]
\item protocol recovery;
\item benchmark success;
\item prediction success;
\item falsification success;
\item physical promotion;
\item physical validation;
\item empirical adequacy;
\item observed-catalog recovery;
\item simulation-only promotion;
\item fit-only calibration;
\item physical nature realization;
\item unified field theory.
\end{itemize}

All constructions below are finite protocol-interface constructions only.

\section{Protocol Records}

\subsection{Finite protocol row}

A Paper 15 protocol row is a finite tuple
\[
P=(i,t,o,r,\theta,\rho,a,b,c),
\]
where:

\begin{itemize}[leftmargin=*]
\item \(i\) is a bounded protocol identifier;
\item \(t\) is a bounded prediction-target label;
\item \(o\) is a bounded observable label;
\item \(r\) is a bounded regime label;
\item \(\theta\) is a bounded threshold label;
\item \(\rho\) is a bounded rejection-condition label;
\item \(a\) is a bounded audit-status descriptor;
\item \(b\) is a Paper 14 compatibility reference;
\item \(c\) is the non-promotion claim boundary.
\end{itemize}

The row closes `PFP-002' exactly when the schema is finite, the labels are
bounded, the row is auditable, Paper 14 compatibility is referenced only as
schema alignment, and the claim boundary contains no success, recovery,
validation, promotion, physical-nature, or unified-field claim.

\subsection{Target, observable, and regime descriptors}

`PFP-003' refines a closed protocol row with finite descriptors for prediction
target, observable, and regime. These descriptors are selection criteria only.
They do not recover an observed catalog and do not import benchmark success,
prediction success, falsification success, validation, empirical adequacy,
physical promotion, or unified-field promotion.

\subsection{Threshold and rejection-condition descriptors}

`PFP-004' adds finite threshold and rejection-condition descriptors. The
threshold is a finite criterion label. The rejection condition is a
predeclared protocol rule. Closing this rung does not mean that any rejection
condition has fired or that any alternative has been falsified.

\subsection{Paper 14 compatibility}

`PFP-005' binds the protocol stack to the frozen Paper 14 commit, formal
endpoint, and final certificate as finite schema alignment:
\[
\texttt{Paper14DiscriminatingBenchmarksTheoremContract.closed}.
\]
This compatibility is not benchmark success and is not prediction or
falsification success.

\subsection{Stability and reproducibility}

`PFP-006' defines finite descriptors for protocol stability, coarse-graining,
and reproducibility. Reproducibility means compatibility of finite records
under preserved labels, audit rows, and claim boundaries. It is not reproduced
empirical success, physical validation, benchmark success, prediction success,
or falsification success.

\subsection{No-hidden-import audit}

`PFP-007' is a fail-closed audit. It closes only when every route that could
silently import recovery, benchmark success, prediction success,
falsification success, validation, empirical adequacy, observed-catalog
recovery, simulation-only promotion, fit-only calibration, physical-nature
realization, or unified-field promotion remains explicitly blocked.

\section{Theorem}

\textbf{Paper 15 Prediction And Falsification Protocols Theorem.}
Assume the closed upstream Paper 14 conditional discriminating benchmarks
certificate and the closed Paper 15 rungs `PFP-001' through `PFP-007'. Then
there exists a finite, network-native prediction-and-falsification protocol
interface whose rows specify bounded protocol identifiers, prediction targets,
observables, regimes, thresholds, rejection conditions, Paper 14 compatibility
references, stability and reproducibility descriptors, and audit statuses.
The interface is closed only as conditional finite protocol structure and
preserves the no-recovery, no-success, no-validation, no-promotion,
no-physical-nature, and no-unified-field claim boundary.

\section{Proof}

The Lean formalization defines contracts for each rung:

\begin{itemize}[leftmargin=*]
\item `PFP001UpstreamBindingContract';
\item `PFP002FiniteProtocolRecordContract';
\item `PFP003PredictionTargetRegimeContract';
\item `PFP004FalsificationThresholdContract';
\item `PFP005Paper14BenchmarkCompatibilityContract';
\item `PFP006StabilityReproducibilityContract';
\item `PFP007NoHiddenPromotionValidationSuccessAuditContract';
\item `PFP008FinalConditionalCertificateContract'.
\end{itemize}

Each contract has a `closed' predicate that is the conjunction of its finite
construction requirements and its claim-boundary guards. For `PFP-001'
through `PFP-007', canonical contracts close their local rung while the Paper
15 theorem contract remains open because the final certificate field is still
false.

The final certificate `PFP-008' consumes the closed canonical contracts for
`PFP-001' through `PFP-007', adds the final conditional-certificate fields,
and preserves every no-claim guard. The formal endpoint is:

\begin{flushleft}
\footnotesize
\texttt{paper15\_pfp008\_final\_conditional\_certificate\_}\\
\texttt{closes\_prediction\_falsification\_protocols\_theorem}
\end{flushleft}

This theorem proves:
\[
\texttt{Paper15PredictionFalsificationProtocolsTheoremContract.closed}
\]
for the final canonical Paper 15 theorem contract.

\section{Audit Surface}

The Rust audit surface mirrors the Lean contracts with finite record types and
regression tests. The final Rust certificate closes only when:

\begin{itemize}[leftmargin=*]
\item the upstream Paper 1--14 binding closes without physical claims;
\item the `PFP-002' through `PFP-007' audit records close;
\item the final certificate is explicitly conditional;
\item all claim boundaries remain false.
\end{itemize}

The regression suite also checks these failure cases:

\begin{itemize}[leftmargin=*]
\item overlong labels, unlinked records, and bad Paper 14 bindings;
\item benchmark-success, prediction-success, and falsification-success
  leakage;
\item empirical-validation leakage and unblocked hidden imports;
\item nonconditional final certificates and unified-field promotion.
\end{itemize}

\section{Reproducibility}

The repository checks are:

\begin{verbatim}
make test-fast
make lean-build
\end{verbatim}

The final state is recorded in `GPD/state.json' with:

\begin{verbatim}
"active_obligation": "NONE"
"prediction_and_falsification_protocols_theorem_closed": true
"paper15_canonical_proof":
  "paper15_pfp008_final_conditional_certificate_"
  "closes_prediction_falsification_protocols_theorem"
\end{verbatim}

\section{Conclusion}

Paper 15 closes the finite prediction-and-falsification protocol interface
theorem conditionally. The result is a theorem about finite protocol records
and their audit boundaries. It is not evidence that predictions succeeded, not
evidence that alternatives were falsified, not a physical validation result,
and not a unified field theory claim.

\end{document}
